% SEGMENT OLP-0027-S01
% Part: sets-functions-relations
% Chapter: size-of-sets-complete

\documentclass[../../../include/open-logic-chapter]{subfiles}

\begin{document}

\olchapter{sfr}{siz}{கணங்களின் அளவு}

\begin{editorial}
இந்த அத்தியாயம் பட்டியலாக்கங்கள், எண்ணிடத்தக்க தன்மை,
எண்ணிடத்தகாத தன்மை ஆகியவற்றை விவாதிக்கிறது.
பல பிரிவுகள் இரண்டு வடிவங்களில் உள்ளன:
பட்டியலாக்கங்களைப் பட்டியல்களாகவோ $\PosInt$ இலிருந்து செல்லும்
மேற்கோர்த்தல் சார்புகளாகவோ கொள்ளும் தொடக்கநிலை வடிவம் ஒன்று;
பட்டியலாக்கங்களை $\Nat$ உடனான இருபுறச் சார்புகளாக வரையறுக்கும்,
மேலும் அருவமான வடிவம் மற்றொன்று.
\end{editorial}

\olimport{introduction}

\olimport{enumerability}

\olimport{zig-zag}

\olimport{pairing}

\olimport{pairing-alt}

\olimport{non-enumerability}

\olimport{reduction}

\olimport{equinumerous-sets}

\olimport{comparing-size}

\olimport{schroder-bernstein}

% SEGMENT OLP-0027-S02
\begin{editorial}
பின்வரும் \olref[sfr][siz][enm-alt]{sec},
\olref[sfr][siz][nen-alt]{sec}, \olref[sfr][siz][red-alt]{sec} ஆகியவை,
\olref[sfr][siz][enm]{sec}, \olref[sfr][siz][nen]{sec},
\olref[sfr][siz][red]{sec} ஆகியவற்றின் மாற்று வடிவங்கள்.
டிம் பட்டன் தனது Open Set Theory நூலில் பயன்படுத்துவதற்காக
இவற்றை உருவாக்கினார்.
இவை சற்றே மேம்பட்ட நிலையில் அமைந்தவை; கணக்கோட்பாட்டுச் சூழலுக்கு
மேலும் பொருத்தமான, பட்டியலாக்கத்திற்கான வேறுபட்ட வரையறையைப்
பயன்படுத்துகின்றன.
அதாவது, பட்டியலிடக்கூடியதாக இருப்பது அல்லது $\PosInt$ இலிருந்து
செல்லும் மேற்கோர்த்தல் சார்பின் வீச்சகமாக இருப்பது என்பதற்குப் பதிலாக,
$\Nat$ உடனோ அதன் ஒரு தொடக்கத் துண்டுடனோ இருபுறச் சார்பு இருப்பது
என்ற வரையறையைப் பயன்படுத்துகின்றன.
\end{editorial}

\olimport{enumerability-alt}
\olimport{non-enumerability-alt}
\olimport{reduction-alt}

\OLEndChapterHook

\end{document}
